\section{Воспроизведение: что способен и не способен выразить исходный дизайн}
\label{sec:repro}

Прежде чем измерять расхождения с исходным исследованием, проверим, что с
его собственными опубликованными числами расхождений нет.

\subsection{Опубликованное преимущество переживает учёт с поправкой на кэш}

Мы пересчитываем итоговую таблицу оригинала по нашей модели стоимости:
берём их опубликованные числа токенов и соединяем их с нашими измеренными
тарифами за токен и нашими измеренными показателями попадания в кэш для
двух архетипичных рантаймов. Результат сужает их преимущество, но не
переворачивает его\live:

\begin{center}
\begin{tabular}{@{}rrr@{}}
\toprule
Горизонт $T$ & отношение токенов & отношение стоимости при учёте кэша \\
\midrule
10  & 1.6$\times$  & \textbf{0.97$\times$} \\
25  & 2.9$\times$  & 1.75$\times$ \\
50  & 5.7$\times$  & 3.42$\times$ \\
100 & 19.0$\times$ & 11.45$\times$ \\
200 & 21.3$\times$ & 12.82$\times$ \\
\bottomrule
\end{tabular}
\end{center}

Из этой таблицы следуют два вывода. Первый: кэширование срезает преимущество
примерно наполовину, но не меняет его знак — внутри исходного дизайна метод
остаётся дешевле, и мы говорим это прямо. Второй: при $T{=}10$ отношение
стоимости уже перешло 1.0, то есть паритет исчезает на самом коротком
горизонте их собственной таблицы, а оригинал этого не замечает, потому что
его метрика такого не видит. Так гипотеза A2 опровергается на одних только
опубликованных данных, без единого нового запуска с нашей стороны.

Эта оценка переносит наши тарифы на их токены, и мы так её и называем:
порядок величины выдерживается на их опубликованных числах, но точные
множители — не измерение того, что они реально потратили.

\subsection{Свойство, которое прячет повторное добывание информации}

В исходном генераторе горизонт $T$ — это число событий среды, и каждый
рантайм отвечает на каждое событие ровно одним обращением к модели\origin.
Число обращений к модели тем самым фиксировано построением и одинаково во
всех руках (во всех сравниваемых условиях).

Это одно свойство решает, что вообще способна увидеть оценка. Главная
стоимость отбрасывания истории — это необходимость заново добывать
выброшенное, а платится она \emph{дополнительными шагами, которые агент
предпринимает по собственной инициативе}. В дизайне, где число шагов —
вход, а не результат измерения, этой стоимости негде проявиться: агент,
забывший факт и вынужденный вывести его заново, всё равно получает ровно
$T$ обращений, и забывание остаётся невидимым в журнале, который ведёт
оригинальная статья.

В агентной постановке число шагов — результат измерения, а не вход. На
нашем корпусе рука «только состояние» (та, что держит лишь суммарное
состояние и выбрасывает транскрипт) тратит \TurnRatioTailZero$\times$ от
числа ходов контроля (той руки, что хранит полный транскрипт) на задачах,
 которые решили \emph{обе} руки\live. Мы измеряем это отношение только на
задачах, где метод справился, поэтому это нижняя граница накладных
расходов, а не среднее с учётом его провалов.

\subsection{Что это подтверждает, а что нет}

Мы воспроизводим механику оригинала, его утверждение о токенах и знак его
преимущества по стоимости с учётом кэша — всё \emph{внутри его собственного
дизайна}. Чего мы не подтверждаем, так это того, что вывод этого дизайна
переносится на агентное выполнение. Обе оценки законны и отвечают на разные
вопросы; но только на один из этих вопросов нужен ответ практику, прежде
чем встраивать архитектуру памяти в кодирующего агента.

Здесь же, а не в сноске, скажем одну оговорку о точности воспроизведения:
наш блок состояния на каждом шаге толще, чем цифра, которую приводит
оригинал\live, поэтому наши абсолютные величины стоимости — верхние
границы. Более тонкий блок усилил бы сокращение промпта и ослабил бы
инверсию стоимости; оба сдвига работают против нашего ключевого результата,
поэтому мы приводим оба.
\note{рука верности исходной статье, которая это изолирует, построена и
предрегистрирована; если она запустится до подачи статьи, её числа заменят
этот абзац}
